\startvolume{Paper IV}{Arithmetic depth, time, and signed coordinates}
\lesson{PAPER IV / 01}{Keep the signs while changing depth}{IV01_signed_seed}
{C_{a,\beta}(k,t)=\sum_{n\ge1}\frac{\mu(n)}{n^{a+it}}y_n^k,\qquad y_n=1-n^{-\beta},\quad a>1,\ \beta>0}
{The seed panel shows signed amplitudes. Squareful integers vanish, while squarefree integers can contribute either sign. The depth panel displays a separate positive filter. It favors larger integers as depth grows. Filtering does not remove the M\"obius signs or replace the signed amplitude with its squared magnitude.}
{Set $a=\beta=2$. The $n=1$ amplitude is one at depth zero. Its filter is zero at every positive depth, so this atom disappears immediately. At $n=2$, the amplitude is $-1/4$ and the depth-five filter is $(3/4)^5$. The complete coefficients begin $C(0,0)\approx0.607927$ and $C(1,0)\approx-0.316011$.}
{Arithmetic decay $a$, scale exponent $\beta$, discrete depth $k$, and logarithmic time $t$ have different jobs.}
{IV Chapter 1, ``One signed seed and two exact evolutions.'' Seed integers through 30; filters through 80. Subsequent principal calculations use $a=\beta=2$ and cutoff $N=200000$.}

\lesson{PAPER IV / 02}{Compute one coefficient in two coordinates}{IV02_finite_difference}
{C_{2,2}(k,0)=\sum_{j=0}^k(-1)^j\binom kj\frac1{\zeta(2+2j)},\qquad |C(k,0)-C_N(k,0)|\le\frac1N}
{The coefficient panel overlays a direct number sum and a finite difference of reciprocal zeta values. Both retain the signed cancellation. The error panel compares their difference with the number-cutoff bound. This bound concerns the omitted infinite tail. It does not describe rounding in a large alternating finite difference, which needs its own precision choice.}
{At depth one, the filter expansion gives $C(1,0)=\zeta(2)^{-1}-\zeta(4)^{-1}$, about $-0.316011$. At depth two it gives $\zeta(2)^{-1}-2\zeta(4)^{-1}+\zeta(6)^{-1}$, about $-0.256997$. Truncating the number sum at $N=200000$ introduces at most $5\times10^{-6}$ absolute error, uniformly in every displayed nonnegative integer depth.}
{Integer depth gives a finite algebraic identity. Infinite number support introduces a separately controlled approximation.}
{IV Chapter 1 finite-difference identity and Worked example 1.1, ``Principal Bernoulli--zeta orbit.'' Depths $0$--$10$; zeta coordinates use 90-digit arithmetic.}

\lesson{PAPER IV / 03}{Watch finite time extract the diagonal}{IV03_finite_time_average}
{\frac1{2T}\int_{-T}^{T}C_N(k,t+\tau)\overline{C_N(\ell,t)}\,dt=\sum_{n,m\le N}a_n\overline{b_m}\operatorname{sinc}\!\bigl(T\log(n/m)\bigr)}
{Each panel follows one component of an exactly integrated finite correlation. Off-diagonal pairs oscillate through their sinc factors. Diagonal pairs retain the lag phase and form the dashed target. Increasing the observation window need not make the error decrease monotonically. Different logarithmic frequency gaps have different cancellation scales.}
{For the pair $n=2$, $m=3$, the off-diagonal factor is $\sin(T\log(2/3))/(T\log(2/3))$. Its first positive zero occurs at $T=\pi/\log(3/2)$, about $7.748$. Other pairs generally remain at that time. The plots keep all $40^2$ pairs, use $k=\ell=3$, and preserve the fixed lag $\tau=0.5$.}
{Finite sinc extraction makes the diagonal mechanism explicit. The infinite time-average theorem has an additional absolute-convergence argument.}
{IV Theorem 2.1, ``Exact time-average identity.'' Here $a_n=\mu(n)n^{-2}y_n^3n^{-i/2}$ and $b_m=\mu(m)m^{-2}y_m^3$; sinc means $\sin x/x$. No infinite-model finite-time bound is inferred.}

\lesson{PAPER IV / 04}{The positive kernel is a complex Gram matrix}{IV04_positive_gram}
{v_{k,t}(n)=\mu(n)n^{-2}y_n^kn^{-it},\qquad K_N(k,\ell;t-s)=\sum_{n\le N}v_{k,t}(n)\overline{v_{\ell,s}(n)}}
{The heatmap shows magnitudes of normalized complex overlaps. The underlying matrix still retains complex phases and is Hermitian. The eigenvalue panel tests the displayed finite packet set. Positivity follows from the feature representation: every quadratic Gram sum is a squared norm. The numerical eigenvalues illustrate that identity rather than proving positivity for every possible packet set.}
{Take the eight packets with depths $0,10,100,1000$ and times zero or one. Divide each finite feature vector by its norm. Every diagonal overlap becomes one. For any coefficients $c_j$, the quadratic form is the norm squared of $\sum_jc_jv_j$, so it cannot be negative, even when individual off-diagonal entries are complex.}
{Time averaging removes arithmetic interference on the diagonal while leaving a positive depth-time geometry.}
{IV Theorem 2.2, ``Positive-definite packet kernel.'' Cutoff $N=200000$; every unnormalized omitted kernel entry is bounded by $1/(3N^3)$. Normalized finite overlaps are displayed.}

\lesson{PAPER IV / 05}{Use Beta to connect depth with Gamma}{IV05_beta_gamma_depth}
{K(k,\ell;\tau)=\frac{\delta}{\beta}B(\alpha+i\tau/\beta,k+\ell+1)+O\!\left((k+\ell)^{-(2a-1/2)/\beta}\right),\quad\alpha=\frac{2a-1}{\beta}}
{At zero lag, multiplying out the leading depth power makes the Gamma limit visible. The Beta term retains a finite-depth correction. The second panel restores a compact interval of nonzero lags and their complex exponent. A single real asymptotic at zero lag would not determine this complete complex law.}
{For $a=\beta=2$, shape is $\alpha=3/2$ and principal squarefree density is $\delta=6/\pi^2$. Writing $m=k+\ell$, the leading zero-lag kernel is $(3/\pi^2)\Gamma(3/2)m^{-3/2}$. At nonzero lag, the power becomes $m^{-3/2-i\tau/2}$. Removing both parts leaves the Gamma factor, whose real part is shown.}
{The continuous shape arises from arithmetic decay and scale. Compact lag intervals are part of the uniform theorem.}
{IV Theorem 4.1, ``Uniform complex packet asymptotic.'' Principal modulus one; $|\tau|\le2$. Finite-cut data illustrate convergence; the theorem's implicit error constant is not fitted from this plot.}

\lesson{PAPER IV / 06}{Remove the carrier before comparing limits}{IV06_carrier_dephasing}
{A_k(\tau)=\frac{K(k,k;\tau)}{K(k,k;0)},\qquad (2k)^{i\tau/\beta}A_k(\tau)\longrightarrow\frac{\Gamma(\alpha+i\tau/\beta)}{\Gamma(\alpha)}}
{The two panels compare real and imaginary components after carrier removal. Without that factor, the logarithmic frequency scale keeps drifting with depth. Normalization removes the changing positive energy; the carrier removes the changing phase reference. Both operations are needed to see the same shape-dependent limit across increasing depths.}
{Equal depth means the combined depth is $2k$. With $\beta=2$, the carrier is $\exp(i\tau\log(2k)/2)$, not $\exp(i\tau\log k/2)$. At zero lag the normalized finite kernel and its limit both equal one. For $k=100,1000,10000$, the figures use this same convention and compare against the shape-$3/2$ Gamma ratio.}
{Changing the phase reference does not alter modulus. It makes complex correlations at different arithmetic scales comparable.}
{IV Corollary 5.1, ``Equal-depth dephasing.'' Compact interval $|\tau|\le4$; finite energies normalize the displayed kernels. Number-cut errors and normalized bounds are recorded separately.}

\lesson{PAPER IV / 07}{An exact envelope can have non-Gaussian tails}{IV07_principal_envelope}
{\left|\frac{\Gamma(3/2+i\tau/2)}{\Gamma(3/2)}\right|^2=\frac{1+\tau^2}{\cosh(\pi\tau/2)}}
{The first panel overlays the Gamma expression with its elementary closed envelope. The second uses a logarithmic vertical scale to compare its tails with a Gaussian having the same limiting variance. Agreement at the initial curvature does not force agreement farther away. The Gamma recurrence and reflection identity determine the full envelope.}
{At lag one the squared modulus is $2/\cosh(\pi/2)$. At zero it is one. The limiting frequency variance is $\psi_1(3/2)/4\approx0.233701$, so the variance-matched Gaussian has squared characteristic modulus $e^{-0.233701\tau^2}$. The exact principal expression instead has an exponential tail multiplied by a quadratic factor.}
{A few moments determine local behavior. A closed characteristic-function formula determines every displayed lag.}
{IV Worked example 5.2, ``Closed principal envelope.'' Exact formulas on $0\le\tau\le8$. The Gaussian is an explicitly variance-matched teaching comparison, not an asserted approximation theorem.}

\lesson{PAPER IV / 08}{Center the logarithmic frequency law}{IV08_centered_log_frequency}
{Z_m=\log n-\beta^{-1}\log m\Longrightarrow-\beta^{-1}\log U,\qquad U\sim\operatorname{Gamma}(\alpha,1)}
{The histograms use positive squarefree packet weights after subtracting the drifting log scale. The limiting density remains continuous even though each finite spectrum consists of integer frequencies. The moment panel compares finite means and variances with their polygamma targets. Histogram convergence alone does not justify convergence of every unbounded moment.}
{For $\alpha=3/2$ and $\beta=2$, the limiting density is $2e^{-3z-e^{-2z}}/\Gamma(3/2)$. Its mean is $-\psi(3/2)/2\approx-0.018244987$ and its variance is $\psi_1(3/2)/4\approx0.23370055$. The finite weights are proportional to $\mu(n)^2n^{-4}(1-n^{-2})^m$, with their full finite mass used for normalization.}
{Log-frequency centering connects the arithmetic active scale with a fixed log-Gamma probability law.}
{IV Corollary 6.1, ``Polygamma cumulants.'' Histograms at $m=1000,10000$; finite moment records through $m=100000$. Preserve the source's moment-convergence proviso when claiming cumulant limits.}

\lesson{PAPER IV / 09}{Compare depth ratios before absolute depths}{IV09_cross_depth_overlap}
{\mathcal O(k,\ell)=\frac{K(k,\ell;0)}{\sqrt{K(k,k;0)K(\ell,\ell;0)}}\longrightarrow\left(\frac{2\sqrt r}{1+r}\right)^\alpha,\qquad r=k/\ell}
{The first panel compares several absolute depth scales against one ratio profile. Equal depths give overlap one. Exchanging the depths leaves the overlap unchanged, so reciprocal ratios share a value. The second follows a single four-to-one ratio while both depths increase. This separates ratio dependence from finite-depth arithmetic fluctuations.}
{Take $\alpha=3/2$ and $k=4\ell$. The limiting overlap is $(4/5)^{3/2}\approx0.715541753$. The same value occurs at ratio $1/4$. At ratio one the limit is exactly one. The displayed comparisons keep ratios between $1/16$ and $16$, and round each sampled larger depth to an integer before evaluating the finite kernel.}
{Comparable depths retain overlap even at large scale. Multiplicative depth separation is measured by a log-depth gap.}
{IV Theorem 7.1, ``Universal scale overlap,'' and Worked example 7.2, ``A four-to-one depth ratio.'' Compact positive ratio domain; finite number-cut sums are not asymptotic proofs.}

\lesson{PAPER IV / 10}{Use the actual character modulus}{IV10_character_sectors}
{\chi_j(2^r)=i^{jr}\quad(\bmod\ 5),\qquad\chi_j(n)=0\ \text{if }5\mid n,\qquad K_{\chi,\psi}=\sum_n\mu(n)^2\chi(n)\overline{\psi(n)}n^{-4-i\tau}y_n^{k+\ell}}
{Each matrix compares all four characters modulo five at one equal depth. Diagonal sectors have unit normalized magnitude. Off-diagonal sectors use nonprincipal character products and weaken with depth. The coprimality restriction is determined by the actual modulus. Replacing it with conductor can change even the principal sector's support and normalization.}
{The units are the cycle $1,2,4,3$ generated by two. For $j=1$, their values are $1,i,-1,-i$, while multiples of five have value zero. The principal character modulo five has conductor one but still vanishes at five. Its squarefree density is $(5/6)/\zeta(2)$, not $1/\zeta(2)$.}
{Distinct characters must share a fixed actual modulus; their product is then nonprincipal and the stated decoherence bound applies.}
{IV Theorem 8.1, ``Character depth-time decoherence,'' with the consolidation's actual-modulus correction. Depths $10,10000$, zero lag; comparable-depth and compact-lag hypotheses govern the asymptotic theorem.}

\lesson{PAPER IV / 11}{Change generating coordinates without losing signs}{IV11_resolvent_borel}
{\begin{gathered}(1-z)\mathcal C_a(z,t)+z\mathcal C_{a+\beta}(z,t)=L(a+it,\chi)^{-1},\\ e^{-x}\sum_{k\ge0}C_a(k,t)\frac{x^k}{k!}=\mathcal R_a(x,t)/x.\end{gathered}}
{The first panel combines two signed resolvents using the parameter-shift identity. The result returns the initial signed coefficient. The second changes ordinary discrete-depth coordinates to normalized exponential coordinates. Removing $e^x$ reveals the Riesz companion. These identities retain the signs throughout; positivity or zero-freeness does not follow from a coordinate change.}
{For the principal choice $a=\beta=2$, $t=0$, the shifted combination returns $1/\zeta(2)$. The normalized Borel value is $\sum_n\mu(n)n^{-2}e^{-x/n^2}$. At $x=0$, define $\mathcal R(x)/x$ by this continuous expression, which returns the same initial coefficient. The finite number cutoff and depth-40 EGF truncation introduce separate errors.}
{Resolvent shifts and Borel--Riesz normalization are exact transformations of one signed family.}
{IV Theorem 9.1, ``Parameter-shift equation,'' and Chapter 10, ``The Borel--Riesz bridge.'' Real $z\in[-0.6,0.8]$, $x\in[0,8]$; bounds for number support and EGF degree are recorded.}

\lesson{PAPER IV / 12}{State the signed Pad\'e contract precisely}{IV12_signed_positive_pade}
{\begin{gathered}\Delta_m\ne0\ \Longrightarrow\ \mathcal C-P/Q=O(z^{2m}),\quad \deg P\le m-1,\ \deg Q\le m,\\ \deg Q=m\ \text{and simple zeros}\ \Longrightarrow\ m\text{ simple packets}.\end{gathered}}
{The node panel contrasts positive and signed effective models. Positive nodes remain in the physical interval; signs can produce a node outside it. The moment panel shades exactly the contracted indices. Later moments are predictions and can differ. Neither numerical nonbreakdown at one order nor this finite contract establishes universal stability.}
{The first signed node is $c_1/c_0\approx-0.519817755$. It already leaves $[0,1)$. Separately, the order-four models use the first eight principal moments evaluated with 90 digits. Their first eight reconstructed moments agree at that precision. Moments eight through eleven have nonzero discrepancies. Full denominator degree and separated simple roots are checked for these displayed models.}
{Hankel nonbreakdown gives rational matching. A simple-packet expansion also needs full denominator degree and simple zeros.}
{IV Theorem 11.1, ``Finite signed Pad\'e reduction,'' with the consolidation correction, and Worked example 11.2, ``The first signed packet leaves the physical interval.'' Numerical examples; plotted errors are floored at $10^{-90}$.}
